% !TeX root = ../../../../main.tex
% sections/chap3/subsections/assumptions/lop.tex

\subsection{一物一価の法則}
\label{sec:chapter3_assumptions_lop}
自国と外国における同一財の価格については、財の国際的な裁定取引により一物一価の法則が成立すると仮定する。

\paragraph{自国財の一物一価}
\label{sec:chapter3_assumptions_lop_home}
自国で生産される財の価格について、自国通貨建て価格\(p_t^h\)と外国通貨建て価格\(p_t^{h*}\)の間には、名目為替レート\(e_t^{/*}\)を介して以下の関係が成立する。
\begin{equation}
p_t^h=e_t^{/*}p_t^{h*}
\label{form:chapter3_assumptions_lop_home_eq}
\end{equation}
もし一物一価の法則が成立せず\(p_t^h<e_t^{/*}p_t^{h*}\)であったとする。
このとき家計は自国において自国財を\(p_t^h\)で買い、外国において\(p_t^{h*}\)で売ることができる。
この\(p_t^{h*}\)を為替市場において自国通貨に戻せば\(e_t^{/*}p_t^{h*}\)となるが、
\(p_t^h<e_t^{/*}p_t^{h*}\)であったのだから\(0<e_t^{/*}p_t^{h*}-p_t^h\)の利益を得ることになる。
一物一価の法則とは、このような裁定取引の結果として名目為替レート\(e_t^{/*}\)や
財の価格\(p_t^h,p_t^{h*}\)が調整され等号が成立すると仮定することにほかならない。

\paragraph{外国財の一物一価}
\label{sec:chapter3_assumptions_lop_foreign}
同様に、外国で生産される財についても以下の等式が成立する。
\begin{equation}
p_t^f=e_t^{/*}p_t^{f*}
\label{form:chapter3_assumptions_lop_foreign_eq}
\end{equation}